\documentclass{situationist-preprint}

\title{Why the List of Goods Is Closed}
\msid{ETH-01-d}
\msfamily{ETH-01}
\msdiscipline{Normative and Applied Ethics}
\mstype{Consequence}
\mskeywords{the good; objective list; completeness; supervenience; method of isolation; value pluralism}

\begin{document}
\maketitle

\begin{abstract}
Objective-list theories of the good are often charged with arbitrariness: the list could always be longer, and nothing in the theory says when to stop. We show that, for the account of the good defended in ETH-01, this charge has a precise answer. Completeness of the list is the thesis that the value of outcomes supervenes on the seven components together. We show that completeness is equivalent to the failure of every further candidate to pass a relative isolation test, that it is stable under further idealization, and that, combined with the plurality theorem of ETH-01-a, it makes the seven components a minimal supervenience base. We then test the principal candidates for an eighth component and show that each fails.
\end{abstract}

\section{Closure as supervenience}

ETH-01 holds that the good is what idealized reflective endorsement converges on (Definition~3.2 of ETH-01-a), and that it has seven irreducible components: well-being, knowledge, achievement, relationships, fairness, beauty and autonomy. ETH-01-a represents an outcome by a profile assigning a level to each component, and proves (Theorem~5.2) that the value $V$ assigned by idealized endorsement does not supervene on any set of components that omits one of the seven. That is the plurality half of the conclusion. The closure half is the claim that nothing needs to be added: the list is complete if and only if every difference in value is accounted for by a difference in some component, which is to say that $V$ supervenes on the seven together. Volume~VI of ETH-01 argues for this thesis at length. Here we derive what closure amounts to and why the usual objection to closed lists does not touch it.

Say that a feature $F$ of outcomes \emph{passes the isolation test relative to} a set $S$ of components if there are two outcomes that agree on every component in $S$, differ in $F$, and are ranked differently by idealized reflective endorsement. This generalizes Definition~4.1 of ETH-01-a, which is the case in which $S$ consists of all components other than $F$, and which is itself a contrastive refinement of Moore's method of isolation \citep[\S 112]{Moore1903}.

\begin{proposition}\label{p:equiv}
$V$ supervenes on a set $S$ of components if and only if no feature of outcomes passes the isolation test relative to $S$.
\end{proposition}

\begin{proof}
If some feature passes relative to $S$, the two outcomes it supplies agree on $S$ and differ in value, so $V$ does not supervene on $S$. Conversely, if $V$ does not supervene on $S$, there are two outcomes that agree on $S$ and differ in value; they differ in some feature, since otherwise they would be the same outcome, and that feature passes relative to $S$.
\end{proof}

\begin{corollary}\label{c:minimal}
If the list is complete, the seven components form a minimal supervenience base for $V$: $V$ supervenes on them, and on no proper subset of them.
\end{corollary}

\begin{proof}
The first half is completeness. The second is Theorem~5.2 of ETH-01-a, since every proper subset omits at least one of the seven.
\end{proof}

Proposition~\ref{p:equiv} shows what the objection from arbitrariness must do. It is not enough to name something else that people value. The objector must produce two outcomes that are alike in well-being, knowledge, achievement, relationships, fairness, beauty and autonomy, and that idealized endorsement nonetheless ranks differently. Corollary~\ref{c:minimal} shows what is at stake if she succeeds: the list would be incomplete, but no component of it would thereby be shown redundant.

\section{Stability under idealization}

A second form of the objection holds that whatever list we arrive at reflects the limits of our present reflection, and that a better-informed judge would find further goods. On the present account this cannot happen in the way the objection supposes. Idealized reflective endorsement is already the attitude of a judge who is fully informed about the outcomes compared and their alternatives, imagines them vividly, reasons without error and gives equal weight to everyone affected (Definition~3.1 of ETH-01-a). The list records what that judge's rankings depend on. A further idealization of the judge is not a change in the standard but a better estimate of it. What may change is our knowledge of which features pass the relative isolation test, and that is a matter of testing candidates, not of waiting for a better judge.

This also explains why the list is short. The judge of Definition~3.1 ranks outcomes from a standpoint that gives weight to everyone affected. An outcome can bear on that standpoint in a limited number of ways: by how the lives of those affected go for them (well-being); by how their minds and efforts are related to the world (knowledge, achievement); by how they are related to one another (relationships); by how goods are distributed among them (fairness); by the extent to which their lives are directed by their own choices (autonomy); and by how the world is apart from any of them (beauty). The seven components are the ways an outcome can matter to such a judge, and a further component would have to be a further way.

\section{Candidates for an eighth component}

We test the candidates most often proposed, using Proposition~\ref{p:equiv}.

\paragraph{Pleasure and contact with reality.} Both are internal to well-being. ETH-05 argues that well-being includes contact with reality; pleasure is a component of well-being on every view. Two outcomes alike in everyone's well-being do not differ in either in a way that idealized endorsement registers.

\paragraph{Virtue.} Virtue is a disposition of agents. Its value for its possessor is a matter of her well-being, of which ETH-11 holds it to be a necessary but not sufficient condition, and its value for others is a matter of their well-being, relationships and the fairness with which they are treated. Hold all of these fixed, and with them every effect of the virtue, and the remaining difference is a difference in dispositions that make no difference to anyone. We find no ranking of such outcomes on reflection, and the tradition that places virtue on the list of goods has typically defended it by appeal to its contribution to a good life, not by such cases \citep{Hurka1993}.

\paragraph{Life.} \citet{Finnis1980} lists life as a basic good. The value of a life's continuing is the value of what it contains for the one who lives it and for others, and two outcomes alike in all seven components do not differ in the value of anyone's life. Life without any of the seven, as in permanent unconsciousness, is not something idealized endorsement prefers for its own sake.

\paragraph{Religion and play.} Finnis also lists religion, understood as harmony with a more-than-human source of meaning, and play. Where the source exists, harmony with it is a matter of knowledge and relationship; where it does not, the outcomes compared do not differ in the relevant fact. Play is valuable for the enjoyment and achievement it contains and drops out when these are held fixed.

\paragraph{Equality and rights.} Equality, so far as it has value beyond the relief of the worse off, is an aspect of fairness. The respect of rights is valuable in outcomes so far as it secures autonomy and fairness; whether rights also constrain action independently of the value of outcomes is a question about the right, not the good, which ETH-02 shows can be represented either way.

No candidate passes the isolation test relative to the seven. This does not prove completeness, since the candidates are not exhausted by those listed here, but it shows that the list is closed against the candidates that have been proposed, and Proposition~\ref{p:equiv} shows that no other form of argument could show it open.

\begin{objection}
Ross included virtue among intrinsic goods, judging that a world containing virtue would be better than one without it even if both contained the same pleasure and knowledge \citep{Ross1930}.
\end{objection}

\begin{reply}
Ross's comparison holds pleasure and knowledge fixed, not well-being, relationships and fairness. Since a world with more virtue in it is, other things equal, a world in which people treat one another better, his comparison does not isolate virtue from the components that virtue affects. Restated with all seven held fixed, it no longer elicits the verdict.
\end{reply}

\section{Conclusion}

The list of goods is closed in a precise sense: completeness is supervenience of value on the seven components together, which is equivalent to the absence of any feature that passes the isolation test relative to them. With the plurality theorem of ETH-01-a, it makes the seven a minimal supervenience base. The charge of arbitrariness is answered by the method that generated the list.

\begin{thebibliography}{9}\small
\bibitem[Finnis(1980)]{Finnis1980} Finnis, J. (1980). \emph{Natural Law and Natural Rights}. Clarendon Press.
\bibitem[Hurka(1993)]{Hurka1993} Hurka, T. (1993). \emph{Perfectionism}. Oxford University Press.
\bibitem[Moore(1903)]{Moore1903} Moore, G.~E. (1903). \emph{Principia Ethica}. Cambridge University Press.
\bibitem[Ross(1930)]{Ross1930} Ross, W.~D. (1930). \emph{The Right and the Good}. Clarendon Press.
\end{thebibliography}

\end{document}
